\subsection{Proof of the original comparison modulus}
\begin{proof}[Proof of Theorem~\ref{thm:modulus120}]
Throughout this proof $d=2$.  Let
$\boldsymbol b=(1/2,1/4)^\top$ and $\boldsymbol s=(1,-1)^\top$.
For a terminal state $y$, normalize its feasible action by
$a=A(y)v$, where $v\in[0,1]$ is held fixed when differentiating
in $y$.  Put
\[
 m(y,v)=F(y,A(y)v,0),\qquad
 k(y,v)=5\{m_1(y,v)-m_2(y,v)\}.
 \]
The original domain argument in
Proposition~\ref{prop:investmentfoundations119} puts $m$, and
every $m+\boldsymbol s z$ used here, in the state cube.
The zero-sum shock leaves the aggregate squared hinge unchanged.
Direct expansion of the original terminal cost therefore gives
\[
 g(m+\boldsymbol s z)=g(m)+kz+5z^2.
\]
By the exact original tail formula in
Theorem~\ref{thm:compiler109}, the normalized last-date objective is
\[
 \widetilde Q(y,v)=S(y)+h(A(y)v)
             +\beta\{g(m(y,v))+\mathcal R_\alpha(k(y,v))\}.
\]
The function $\mathcal R_\alpha$ is continuously differentiable and
\[
 \mathcal R_\alpha'(k)
   =\operatorname{clip}\left(\frac{\alpha k}{10t},
                                      -\alpha R,\alpha R\right).
 \]
Its derivative has Lipschitz constant $\alpha/(10t)$, including
the two clipping junctions.

We establish a uniform upper state-curvature bound for
$\widetilde Q(\cdot,v)$.  For $d=2$, the target, cyclic, and
squared-hinge terms in $S$ have Hessian operator norms bounded by
$2$, $1$, and $8$, respectively.  Thus
$\nabla^2S\preceq11I$ wherever the second derivative is defined.
The same decomposition for $g$ gives
$0\preceq\nabla^2g\preceq13I$.
Its gradient coordinates satisfy
\[
 |\partial_jg(y)|
 \le 4|y_j-5/8|+\tfrac12|y_j-y_{j+1}|
                  +4(1/2-y_1-y_2)_+
 \le5.
\]
The state Jacobian of $m$ is nonnegative on the cube.  Its row
sums are at most $3/4$, and its column sums at most $47/64$.
Indeed the unnormalized transition has both sums at most $11/16$;
the term $\boldsymbol b v\nabla A^\top$ adds at most $1/16$
to a row and $3/64$ to a column.  Hence
\[
 \|D_y m\|_2\le\sqrt{(3/4)(47/64)}\le3/4.
\]
The only second state derivatives of $m_j$ are the mixed
derivatives $-1/16$ of $-y_1y_2/16$; the normalized capacity term
is affine.  In a state direction $w$,
\[
 \left|\sum_{j=1}^2\partial_jg(m)\,
                         D_y^2m_j[w,w]\right|
 \le \frac{10}{8}|w_1w_2|
 \le\frac58\|w\|_2^2.
\]
The chain rule consequently gives the uniform upper bound
\[
 D_y^2(g\circ m)[w,w]
 \le\left\{13(3/4)^2+5/8\right\}\|w\|_2^2
 =\frac{127}{16}\|w\|_2^2.
\]
Also $h''(a)=2+48a^2\le5$ on $[0,1/4]$ and
$\|\nabla A\|_2^2=1/128$.  The stage upper bound is therefore
$11+5/128=1413/128$.

For the risk-motion term, the bilinear transition terms cancel
in the difference of the two coordinates:
\[
 k(y,v)=\frac{35}{16}(y_1-y_2)+\frac54A(y)v.
\]
This is affine in $y$ for fixed $v$, and
\begin{align*}
 \nabla_y k&=(35/16+5v/64,\ -35/16+5v/64)^\top,\\
 \|\nabla_y k\|_2^2
 &=2(35/16)^2+2(5v/64)^2
 \le\frac{19625}{2048}.
\end{align*}
It follows that the upper Hessian contribution of
$\mathcal R_\alpha(k(y,v))$ is at most
$19625\alpha/(20480t)$.  Combining these estimates gives
$D_y^2\widetilde Q[w,w]\le M_\alpha\|w\|_2^2$.
Squared hinges and the tail formula have matching first derivatives
at their junctions.  Their restrictions to a line have absolutely
continuous first derivatives, so integrating this almost-everywhere
second-derivative bound along a segment proves that
\[
 y\longmapsto\widetilde Q(y,v)-\tfrac12M_\alpha\|y\|_2^2
\]
is concave on the cube, uniformly over $v\in[0,1]$.

An infimum of concave functions is concave: its value at a convex
combination is at least the same combination of the two infima.
Since the normalized feasible interval is fixed, this proves
semiconcavity of
$V_1(y)-M_\alpha\|y\|_2^2/2$.
The retained original joint curvature and the affine capacity graph
give convexity and a unique continuous minimizing $v(y)$.
For clarity, continuous differentiability does not require
differentiability of that minimizer.  Evaluate a value increment at
the old minimizer for one inequality and at the new minimizer for
the opposite inequality.  Uniform first-order expansion of the
$C^1$ objective on the compact normalized interval and continuity of
$v(y)$ give
$\nabla V_1(y)=\partial_y\widetilde Q(y,v(y))$ in the interior.
This formula extends continuously to the boundary, and incorporates
the capacity term already displayed in
\eqref{eq:constrainedenvelope119}.

Convexity together with $M_\alpha$-semiconcavity now implies the
claimed gradient modulus.  One explicit justification is to mollify
$V_1$ on any smaller interior cube.  Convolution preserves convexity
and semiconcavity, so the smooth Hessian has every eigenvalue between
$0$ and $M_\alpha$.  Integrating the Hessian along a segment bounds
the difference of the mollified gradients by $M_\alpha$ times its
length.  The gradients converge locally uniformly because $V_1$ is
$C^1$.  Pass to the limit and then use continuity at the boundary.

Finally the two successor points at a fixed innovation differ by
$\boldsymbol b(b-a)$.  Differentiation in the innovation gives
\begin{align*}
 \left|\frac{d}{dz}\{f_b(z)-f_a(z)\}\right|
 &\le\|\boldsymbol s\|_2
    \|\nabla V_1(F(x,b,z))-\nabla V_1(F(x,a,z))\|_2\\
 &\le M_\alpha\|\boldsymbol s\|_2\|\boldsymbol b\|_2|b-a|\\
 &=\frac{\sqrt{10}}4M_\alpha|b-a|
 \le\frac45M_\alpha|b-a|.
\end{align*}
All points and joining segments are in the invariant convex cube.
Integration in $z$ proves \eqref{eq:mixedmodulus120}.
\end{proof}

